% =====================================================================
%  Walkability of the Sudirman Corridor (ENGLISH, IEEE conference template)
%  Companion to main_id.tex (Bahasa Indonesia, report class).
%  Compile on Overleaf: pdfLaTeX + Biber. Shares refs.bib.
%  Research-stage paper: NO field findings until data is collected.
%  Depth target: at least as detailed as main_id.tex (no page limit).
% =====================================================================
\documentclass[conference]{IEEEtran}

\usepackage[T1]{fontenc}
\usepackage[utf8]{inputenc}
\usepackage{graphicx}
\usepackage{amsmath}
\usepackage{textcomp}
\usepackage{xurl}
\usepackage[hidelinks]{hyperref}
\usepackage{csquotes}
\usepackage[backend=biber,style=ieee]{biblatex}
\addbibresource{refs.bib}

\begin{document}

\title{Toward a Comfort-Aware Assessment of Walkability in the Sudirman Corridor, Jakarta:\\
A Mixed-Methods and Computer-Vision Research Design}

\author{%
  \IEEEauthorblockN{Gading Aditya Perdana, Kenrich Heavenly Sandria, Naufal Ammar Rauf,
  Jatayu Muhammad Wicaksono, Sintani Gina Lestari}
  \IEEEauthorblockA{Tim Riset Pijak (Independent Research Team)\\ Jakarta, Indonesia}%
  \thanks{Independent team research conducted during the Apple Developer Academy @ BINUS,
  Cohort 9. Intellectual property belongs to the team. Working challenge statement; results
  follow fieldwork.}%
}

\maketitle

\begin{abstract}
Walking is the most basic form of urban mobility, yet in Greater Jakarta it accounts for under
two percent of commute trips and sits within a city where sidewalks cover a small share of the
road network. The Jalan Jenderal Sudirman corridor is a useful case: a 2017--2022 sidewalk
revitalization pilot with a sharp contrast between a wide main footway and the near-absent
sidewalks one block behind it. This paper presents the research design for a comfort-focused
walkability study of that corridor and its back-streets. We combine a standards-referenced
segment audit, time-windowed pedestrian counts, neutral street intercepts, and an
obstacle-oriented expert interview with a computational layer: computer-vision mapping of the
sidewalk network and streetscape features, a shade and thermal-comfort layer, and interpretable
machine learning for route and mode choice. Weather is logged per observation window from a
single source to test whether conditions change comfort and behavior. The core design is a
paired observation of the main corridor against its back-street to measure the quality gradient
that the prior literature leaves unmeasured. We report the instruments, data sources, benchmarks,
the computational pipeline, and the cross-validation strategy; field results are deliberately
withheld until data collection. The team was flagged for confirmation bias, so the design keeps
questions neutral, separates assumptions explicitly, and reports both presence and absence of
effects.
\end{abstract}

\begin{IEEEkeywords}
walkability, pedestrian environment, thermal comfort, computer vision, machine learning,
mixed methods, transit-oriented development, Jakarta
\end{IEEEkeywords}

\section{Introduction}
Population mobility in Jakarta depends on mass transit (Transjakarta, Commuter Line, MRT), yet
the first and last leg of most trips is made on foot, by motorcycle taxi (ojek/ojol), or by
private vehicle. Walking holds a small share of commute modes: under two percent of trips in the
2023 Greater Jakarta commuter survey are made on foot or by bicycle, with private vehicles
dominant \autocite{goodstats2024}. City-wide sidewalk provision is thin: about 540~km of sidewalk
against 6{,}956~km of road in 2015 \autocite{itdp2019}. Walking in a transit-dense corridor is
therefore best understood as a last-mile connector rather than a primary mode.

The Jalan Jenderal Sudirman corridor is a strong case. It was a flagship of the 2017--2022
sidewalk revitalization program \autocite{beritajakarta2019,itdp2019}, and it shows a sharp
contrast between a wide, engineered main footway and the degraded streets one block behind it. We
call this contrast the gradient, and it is the central object of study.

\subsection{Essential Question and Challenge Statement}
The essential question framing the project is: what determines whether walking in the city feels
comfortable, safe, and dignified, and for whom is that experience realized? The working challenge
statement, kept temporary and to be refined after fieldwork, is to improve the comfort of the
walking experience in the Sudirman transit corridor. The word ``comfort'' is treated as an
assumption to be tested, not a conclusion.

\subsection{Research Questions}
The questions are neutral and two-sided, and they carry no embedded assumption:
\begin{enumerate}
  \item How do pedestrians move through the Sudirman area?
  \item Do pedestrians use the available sidewalks, and what influences this?
  \item What makes walking here feel comfortable, and what makes it uncomfortable?
  \item How do weather (heat, rain) and the availability of shade affect when, where, and how
        people walk, and how do they cope (umbrella, sheltering, switching to ojol)?
  \item How do pedestrians know, or fail to know, route conditions before walking?
  \item How does route choice differ by time of day and by gender?
  \item What obstacles blocked the delivery of comfortable pedestrian facilities in 2017--2022?
        (for the expert interview)
\end{enumerate}

\subsection{Scope and Stance}
The scope is locked to the Sudirman corridor and its immediate back-streets, named by street
rather than by district or by a list of stations. Dukuh Atas is not in the core scope and is
included only where feasible in the field. The team was flagged for confirmation bias, so the
questions stay neutral, assumptions are made explicit and kept separate from the questions, and
findings are withheld until field data is collected. The empty results chapter is intentional.

\section{Related Work}
This review is organized by theme and distinguishes corridor-specific evidence from city or
national context. Where possible each figure is traced to its source, and limitations are stated
openly.

\subsection{Walkability Concepts and Measurement}
Pedestrian-path quality is commonly scored with structured indices. The Pedestrian Environmental
Quality Index (PEQI) rates completeness, safety, and comfort from both physical condition and
perception \autocite{dwisadana2024}. The Global Walkability Index (GWI) uses nine parameters
\autocite{ridhani2015,mulyadi2022}, with categories of highly walkable (index $\geq 70$), waiting
to walk ($50$ to $70$), and not walkable ($\leq 50$) \autocite{mulyadi2022}. At the policy level,
design guidance derives from a Complete Streets framework with completeness, safety, comfort, and
humaneness as its elements \autocite{itdp2019}. Because indices differ in method, scores from
different studies are not directly comparable, so change over time must be read with care.

\subsection{Corridor Condition and User Satisfaction}
On the southern segment of the corridor, a post-revitalization PEQI study scored physical
condition $64.28$ and perception $84.70$, for a combined $74.49$ classified as acceptable
\autocite{dwisadana2024}. The same study lists concrete shortfalls: motorcycles and street
vendors on the footway despite bollards, limited shade, missing pedestrian signals at some
crossings, noise, and a transit distance felt as far. The authors attribute part of the gap
between perception and measured condition to method limits and to users tolerating a footway as
long as it carries their movement; a high perception score therefore does not by itself mean high
comfort. A road-diet revitalization raised GWI walkability by up to $38.98\%$ and was associated
with a $15.41\%$ rise in public-transport users on a corridor bus route, from a 2015 baseline near
$64$ (waiting to walk) on Sudirman and $71$ (walkable) on Thamrin
\autocite{mulyadi2022,ridhani2015}. Because PEQI (2024) and GWI (2015 and 2018) differ in method,
segment, and year, the pattern of improvement followed by residual problems is a hypothesis to
test in the field, not a settled trend. A user-satisfaction study on the same corridor evaluated
paths, street furniture, and signage from the pedestrian side, complementing the index studies
with a user-side view \autocite{apritasari2020satisfaction}.

\subsection{Transit and Last-Mile Walking}
Pedestrian activity concentrates at transit nodes \autocite{dwisadana2024}; at Dukuh Atas, walking
frequency is driven by income, trip purpose, and settlement density \autocite{rakhmatulloh2020},
and a recent study reassesses walkability on a 400~m core radius after the MRT and
transit-oriented development \autocite{napitupulu2025tod}. For the MRT corridor, a
stated-preference discrete-choice study reports Jakarta-specific acceptable walking distances of
about $629$~m for men and $593$~m for women, and about $689$~m near offices versus $547$~m near
residences \autocite{afkara2020walking}. Its own future work asks that pedestrian-pathway
characteristics be modeled as a factor and that richer (multinomial or machine-learning) models be
used, an opening this project takes up. Field reporting documents last-mile friction near a
Sudirman station, where commuters crowd the kerb for ojol while motorcycles line the footway and
pedestrians disperse on foot \autocite{kompas2025ojol}; the two named pedestrians there are young
women, which usefully counterweights the male-commuter profile in the survey studies.

\subsection{Safety, Gender, and Public Space}
Public streets are the most common offline site of harassment nationally, with women reported far
more vulnerable \autocite{krpa2022}, and women's daily mobility is framed as linked trip chains
that demand safe and comfortable space \autocite{udk_mobilitasperempuan}. National crash data
record $10{,}428$ pedestrian victims in 2023, with $54.84\%$ linked to crossing at undesignated
places \autocite{korlantas2023}; this makes the crossing audit (spacing, signals, jaywalking)
relevant to real risk. All of this evidence is context: none is broken down for the Sudirman
corridor specifically.

\subsection{Microclimate, Shade, and Thermal Comfort}
Heat shapes comfort in this tropical setting. Sky View Factor, the openness of the sky above a
footway, relates to pedestrian thermal comfort on a Jakarta sidewalk \autocite{arif2021skyview},
and shade is a recurring shortfall in corridor audits, where the guidance recommends shade roughly
every $300$~m and canopy at crowd points \autocite{dwisadana2024}. Tropical thermal-comfort work
commonly uses a Temperature-Humidity index, with values above roughly $27\,^{\circ}$C treated as
uncomfortable (a threshold this project verifies against its own readings).

\subsection{Informal Use, Policy, and Standards}
Pedestrian space is occupied informally; Jakarta had 324 pedestrian bridges in 2015, mostly near
transit, with leftover space colonized for unplanned uses \autocite{putri2020residual}. This frames
informal occupation as something to record neutrally as a finding, not only to condemn as an
obstruction. A 2018 policy reserved the Sudirman--Thamrin footway for pedestrians and planned to
relocate vendors to side streets behind office buildings \autocite{merdeka2018pkl}, a plausible
driver of the gradient. Vehicle encroachment on sidewalks remains a current city-wide concern
\autocite{kompas2026trotoar}, consistent with the corridor-level finding on motorcycles and
vendors \autocite{dwisadana2024}. National technical standards set a minimum footway width of
$1.85$~m, crossing spacing of $100$ to $200$~m, a maximum ramp slope of $8\%$ (1:12), a minimum
overpass clearance of $5.1$~m, and a protected-crossing threshold above $450$ pedestrians per hour
per effective meter \autocite{binamarga2023}.

\subsection{Computational Methods}
Recent work measures walkability from imagery with deep learning, segmenting street-view panoramas
into elements that become walkability features \autocite{walkabilitydl2024}, with crowdsourced
accessibility labeling schemes for sidewalks \autocite{projectsidewalk2019}. Open-source semantic
segmentation of aerial tiles maps sidewalk, crosswalk, and footpath into a routable network,
allowing sidewalk presence and continuity to be measured at scale \autocite{tile2net2023}.
Shade-aware routing models a sun-avoidance preference and a CoolWalkability metric from building
footprints, heights, and sun position \autocite{coolwalks2025}. Interpretable machine learning,
using gradient boosting with SHAP, extends discrete-choice models of pedestrian route choice
\autocite{routechoiceml2024,afkara2020walking}.

\subsection{Research Gaps}
Three gaps define the space for fieldwork. First, no collected source measures the footway
condition of the back-streets one block behind the corridor; every corridor study assesses the
main axis or a transit node \autocite{dwisadana2024,rakhmatulloh2020,ridhani2015,mulyadi2022}, so
the gradient is implied by city-scale ratios but never measured at street level. Second, the
meaning of ``comfort'' to corridor pedestrians has not been gathered in their own words; sources
score predefined parameters rather than user definitions. Third, the obstacles that made
comfortable facilities hard to deliver are undocumented; policy sources describe what was built,
not what blocked it \autocite{beritajakarta2019,itdp2019}, which is the purpose of the expert
interview. Corridor-specific safety data and day-versus-night, non-commuter coverage are likewise
absent.

\section{Methodology}
\subsection{Research Design}
We use a convergent mixed-methods design that combines three angles of evidence: objective
(standards-referenced audit, counts, sensors, computer vision), observed (pedestrian behavior),
and stated (interviews and validated perception scales). The core is a paired observation of a
main-corridor segment against its back-street to quantify the gradient that the literature leaves
unmeasured.

\subsection{Scope and Units}
The unit of analysis is a segment (one block per side, named by street). Each main-corridor
segment is paired with the back-street block behind it, observed on matched days and times.
Windows span the morning peak (06:00--09:00), midday heat (11:00--14:00), evening peak
(15:00--18:30), after dark, and weekend or Car Free Day. Each window records its weather. Priority
nodes are Dukuh Atas and Sudirman station, GBK--FX, and Ratu Plaza and Menara Mandiri.

\subsection{Instruments}
\begin{itemize}
  \item \textbf{Segment audit.} Physical parameters scored 0 to 2 (sidewalk presence, effective
        width, surface, obstruction, ramp and tactile guiding block, lighting, crossing access,
        shade, drainage), with a GPS pin and photo references.
  \item \textbf{Tally sheet.} Five-minute counts per window (sidewalk versus road, direction to or
        from transit, apparent group and gender), an obstruction inventory, and a weather and
        coping block (umbrella use for sun or rain, shade-seeking, sheltering, switching to ojol).
  \item \textbf{Intercept script.} A street intercept of at most five minutes, consent first, with
        open and two-sided questions.
  \item \textbf{Interview guide and bank.} Neutral per-persona prompts and an obstacle-oriented
        expert interview.
  \item \textbf{Validated perception scales.} A seven-point thermal sensation vote (ASHRAE), a
        perceived safety and comfort scale, and selected Neighborhood Environment Walkability Scale
        items, so ratings are comparable across people and against the literature.
\end{itemize}

\subsection{Sampling}
Sampling is non-probability with quotas crossing persona, location, and time. Personas include
office commuters, women pedestrians (oversampled), students, transit last-mile users, informal
actors (vendors and ojol drivers), back-street residents, and people with reduced mobility.
Recruitment occurs at natural dwell points such as ojol waiting spots, crossings, bus stops, and
warung. The convenience nature of the sample is stated openly; representativeness is not claimed.

\subsection{Data Sources}
\begin{itemize}
  \item \textbf{Weather.} Apple WeatherKit is the primary source (temperature, apparent
        temperature, humidity, precipitation probability, UV, condition), logged per window on
        iPhone or Apple Watch so research and any later application share one source; Open-Meteo is
        a historical fallback and the national agency (BMKG) an official cross-check, with
        attribution shown on outputs.
  \item \textbf{Air quality.} Used as a comfort and health factor.
  \item \textbf{Geometry.} OpenStreetMap for sidewalks, crossings, building footprints and heights,
        with an open digital elevation model.
  \item \textbf{Imagery.} Google Street View and aerial tiles as computer-vision inputs, with
        capture dates recorded.
  \item \textbf{Sensing walk.} A logger for temperature, humidity, and illuminance with GPS,
        carried along each segment and tagged to the segment so readings join the audit.
  \item \textbf{Sentiment.} Geotagged review and social-media mining as a public-discourse layer.
\end{itemize}

\subsection{Benchmarks}
Audit parameters reference the national standard \autocite{binamarga2023}: minimum footway width
$1.85$~m, crossing spacing $100$ to $200$~m, maximum ramp slope $8\%$ (1:12), minimum overpass
clearance $5.1$~m, and the pedestrian-flow unit of persons per effective meter per minute, with a
protected crossing warranted above $450$ pedestrians per hour per effective meter. Walkability uses
PEQI and GWI \autocite{dwisadana2024,mulyadi2022}; acceptable walking distance uses Jakarta-specific
values \autocite{afkara2020walking}; safety context uses national crash data
\autocite{korlantas2023}; thermal comfort references Sky View Factor and a Temperature-Humidity
index \autocite{arif2021skyview}.

\subsection{Computational Methods}
We apply the team's computer-vision and machine-learning skills to widen coverage and objectivity.
\begin{enumerate}
  \item \textbf{Automated sidewalk mapping.} Semantic segmentation of aerial tiles produces
        sidewalk and crossing polygons and a routable centerline network, measuring the gradient at
        scale \autocite{tile2net2023}. The model is trained on other cities, so it requires
        validation against the Jakarta field audit.
  \item \textbf{Streetscape features.} Street-view segmentation yields per-point proportions of
        sidewalk, greenery, sky, and obstructions \autocite{walkabilitydl2024}, with an
        accessibility label scheme \autocite{projectsidewalk2019}.
  \item \textbf{Shade and thermal layer.} Shade is computed from building footprints and heights
        and sun position following the CoolWalks approach \autocite{coolwalks2025}, complementing
        field thermal readings.
  \item \textbf{Behavior modeling.} A discrete-choice baseline \autocite{afkara2020walking} is
        extended with an interpretable model (gradient boosting with SHAP) using
        computer-vision-derived and field features \autocite{routechoiceml2024}.
\end{enumerate}
Each model is treated as a hypothesis generator; corridor-level conclusions come from field data
and ground truth, not from a model alone.

\subsection{Cross-Validation}
Every claim is tested from the three angles. Computer-vision outputs are ground-truthed against the
field audit (for example, segmented width against the measured width and the $1.85$~m standard);
sensor readings are compared with WeatherKit; and audits are scored by multiple raters to test
inter-rater reliability. Agreement across angles strengthens a finding; disagreement is reported as
insight rather than error.

\subsection{Ethics and Neutrality Discipline}
Every intercept begins with verbal consent, with separate consent for any recording. Sensitive and
after-dark intercepts are led by a female team member with a buddy present. Faces in photographs
are blurred, and service terms are respected (WeatherKit attribution, Google terms, OpenStreetMap
license). All questions stay neutral and two-sided, assumptions are kept explicit and separate, and
findings are not written before data is collected. The team records both sidewalk and road use,
both comfortable and uncomfortable conditions, and both high and low counts, because the absence of
an effect is also data.

\section{Expected Results and Discussion}
Field results are withheld until data collection. The analysis will report the measured gradient
between the main corridor and its back-streets, the relationship between weather and observed
comfort behavior, the agreement between computer-vision estimates and the field audit, and the
features that an interpretable model weights for route and mode choice. No conclusion about
corridor comfort or safety is drawn in advance, and the team will report effects that are absent
as carefully as effects that are present.

\section{Conclusion}
This paper sets out a neutral, standards-referenced, and computationally extended research design
for a comfort-focused walkability study of the Sudirman corridor and its back-streets. The
contribution at this stage is the design itself: a paired-gradient measurement, a weather-joined
observation protocol, validated instruments, and a computer-vision and machine-learning pipeline
that is validated against field ground truth and built to close the gaps the literature leaves.
Findings will follow the fieldwork.

\printbibliography

\end{document}
